\documentclass[border=4pt]{standalone}
\usepackage{amsmath,amssymb,mathtools,xcolor,varwidth}
\definecolor{widthorange}{HTML}{E8770A}
\definecolor{boundred}{HTML}{D81B60}
\definecolor{stepblue}{HTML}{1E6FE0}
\definecolor{basegreen}{HTML}{00A35A}
\begin{document}
\begin{varwidth}{28em}
\large\bfseries
Let \textcolor{widthorange}{$AF$} be taken equal to the greatest of the breadths, and complete \textcolor{boundred}{the parallelogram $FAaf$} upon it. Because this single rectangle is at least as tall and as wide as any of the differences between the circumscribed and inscribed steps, \textcolor{boundred}{$FAaf$} is greater than the whole excess of the circumscribed figure over the inscribed one. But as the bases are all diminished \emph{in infinitum}, \textcolor{widthorange}{the greatest breadth $AF$} likewise shrinks toward nothing, so \textcolor{boundred}{$FAaf$} becomes less than any given rectangle. Hence the trapped error, lying wholly within \textcolor{boundred}{$FAaf$}, also vanishes, and \textcolor{stepblue}{the inscribed} and circumscribed figures along \textcolor{basegreen}{the base $AE$} become ultimately equal exactly as in \textbf{Lemma~II}. \textit{Q.E.D.}
\end{varwidth}
\end{document}
